%% 견본 장: 디자인 확인용. 내용은 제2부 「부동소수점 수」 장의 앞부분 초고.
\part{계산하는 기계}

\motto{컴퓨터는 실수를 계산하지 않는다. 실수에 가장 가까운, 셀 수 있는 몇 개의 수를 계산할 뿐이다.}
\chapter{부동소수점 수}
\label{ch:float}

\abstract{수학에서 실수는 빈틈없이 이어진 직선이지만, 컴퓨터의 메모리는 유한하다. 64비트로 적을 수 있는 서로 다른 값은 많아야 $2^{64}$개이므로, 컴퓨터는 실수 전체 가운데 아주 작은 부분집합만을 정확히 담고 나머지는 그 부분집합의 원소 가운데 하나로 \emph{반올림}해서 다룰 수밖에 없다. 어떤 부분집합을 고를 것인가, 반올림은 어떻게 할 것인가, 그리고 그렇게 생긴 오차는 계산을 거치며 어떻게 자라는가. 이 장에서는 오늘날 거의 모든 컴퓨터가 따르는 IEEE~754 표준이 앞의 두 질문에 내놓은 답을 살피고, 셋째 질문에 답하기 위한 가장 기본적인 도구인 반올림 모형을 세운다.}

\section{왜 부동소수점인가}
\label{sec:why-float}

정해진 비트 수로 실수를 적는 가장 단순한 방법은 소수점의 자리를 고정하는 것이다. 예컨대 32비트 가운데 앞 16비트를 정수 부분에, 뒤 16비트를 소수 부분에 쓰기로 약속하면 $2^{-16}$ 간격으로 놓인 수들을 정확히 적을 수 있다. 이를 \term{고정소수점}{fixed point} 표현이라 한다. 고정소수점은 덧셈과 뺄셈이 정수 연산과 똑같다는 장점이 있지만, 적을 수 있는 수의 \emph{크기}가 좁다는 치명적인 약점이 있다. 위의 약속으로는 $2^{15}$보다 큰 수를 적을 수 없고, $2^{-16}$보다 작은 수는 모두 0이 된다. 과학 계산에서는 행성의 질량($\approx 10^{24}$\,kg)과 전자의 질량($\approx 10^{-30}$\,kg)이 한 식에 나란히 나타나는 일이 드물지 않으므로, 이런 표현으로는 버틸 수 없다.

부동소수점은 이 문제를 우리가 과학적 표기법에서 이미 쓰고 있는 방법으로 푼다. $6.022\times10^{23}$처럼 유효숫자 몇 자리와 지수를 따로 적으면, 유효숫자의 자릿수는 고정한 채 지수만 바꾸어 아주 큰 수와 아주 작은 수를 같은 \emph{상대} 정밀도로 적을 수 있다. 소수점이 지수에 따라 ``떠다닌다''는 뜻에서 \term{부동소수점}{floating point}이라는 이름이 붙었다.

\begin{definition}[부동소수점 수의 집합]
\label{def:float-set}
기수 $\beta\ge2$, 정밀도 $p\ge1$, 지수의 범위 $e_{\min}\le e_{\max}$가 주어졌을 때
\begin{equation}
\label{eq:float-set}
\mathbb{F} = \bigl\{\, \pm\, m\,\beta^{\,e-p+1} \;:\; m\in\mathbb{Z},\ 0\le m<\beta^{p},\ e_{\min}\le e\le e_{\max} \,\bigr\}
\end{equation}
를 \term{부동소수점 수의 집합}{set of floating-point numbers}이라 하고, $m$을 \term{가수}{significand}, $e$를 \term{지수}{exponent}라 한다. 특히 $\beta^{p-1}\le m$인 원소를 \term{정규수}{normal number}, $m<\beta^{p-1}$이고 $e=e_{\min}$인 원소를 \term{비정규수}{subnormal number}라 한다.
\end{definition}

\cref{def:float-set}에서 가수 $m$을 정수로 둔 것은 증명을 깔끔하게 하려는 선택일 뿐이다. 흔히 보는 $1.d_1d_2\ldots d_{p-1}\times\beta^{e}$ 꼴은 $m=\beta^{p-1}+d_1\beta^{p-2}+\cdots+d_{p-1}$로 두면 \eqref{eq:float-set}의 원소와 정확히 같다. 이 책에서 가장 자주 만날 형식의 매개변수를 \cref{tab:ieee-params}에 모았다.

\begin{table}[t]
\begin{sdtable}{colspec={lccccc}, row{1}={font=\small\bfseries}}
\toprule
형식 & 비트 & $p$ & $e_{\min}$ & $e_{\max}$ & $u=2^{-p}$ \\
\midrule
binary32 (단정밀도) & 32 & 24 & $-126$ & $127$ & $2^{-24}\approx5.96\times10^{-8}$ \\
binary64 (배정밀도) & 64 & 53 & $-1022$ & $1023$ & $2^{-53}\approx1.11\times10^{-16}$ \\
\bottomrule
\end{sdtable}
\caption{IEEE~754의 두 이진 형식. $u$는 \cref{sec:rounding}에서 정의할 반올림 단위이다.}
\label{tab:ieee-params}
\end{table}

\subsection{수직선 위의 부동소수점 수}

$\mathbb{F}$의 원소는 수직선 위에 고르게 놓여 있지 않다. 지수가 $e$인 정규수는 구간 $[\beta^{e},\beta^{e+1})$에 간격 $\beta^{e-p+1}$로 놓이므로, 지수가 하나 커질 때마다 간격이 $\beta$배로 벌어진다. \cref{fig:float-line}은 $\beta=2$, $p=3$인 장난감 형식에서 이를 그린 것이다.

\begin{figure}[t]
\centering
\begin{tikzpicture}[x=13mm]
  \draw[-{Latex[length=2mm]}, black] (-0.1,0) -- (8.4,0);
  % 비정규수와 정규수: beta=2, p=3, e_min=-1
  \foreach \x in {0,0.125,0.25,0.375}{\draw[black!45, line width=0.8pt] (\x,-0.1)--(\x,0.1);}
  \foreach \x in {0.5,0.625,0.75,0.875,1,1.25,1.5,1.75,2,2.5,3,3.5,4,5,6,7}{\draw[black, line width=0.8pt] (\x,-0.12)--(\x,0.12);}
  \foreach \x in {0,1,2,4}{\node[below=3pt, font=\footnotesize] at (\x,0) {$\x$};}
  \node[below=3pt, font=\footnotesize] at (0.5,0) {$\tfrac12$};
  \node[below=3pt, font=\footnotesize] at (7,0) {$7$};
  \draw[decorate, decoration={brace, amplitude=4pt}, black] (4,0.25) -- (7,0.25)
     node[midway, above=5pt, font=\footnotesize, text=black] {간격 $1$};
  \draw[decorate, decoration={brace, amplitude=4pt}, black] (2,0.25) -- (3.5,0.25)
     node[midway, above=5pt, font=\footnotesize, text=black] {간격 $1/2$};
  \draw[decorate, decoration={brace, amplitude=3pt, mirror}, black] (0,-0.55) -- (0.375,-0.55)
     node[midway, below=4pt, font=\footnotesize, text=black] {비정규수};
\end{tikzpicture}
\caption{$\beta=2$, $p=3$, $e_{\min}=-1$, $e_{\max}=2$인 형식의 음이 아닌 원소. 지수가 하나 커질 때마다 간격이 두 배가 된다. $0$과 가장 작은 정규수 $1/2$ 사이는 비정규수가 같은 간격으로 메운다.}
\label{fig:float-line}
\end{figure}

\section{반올림과 그 오차}
\label{sec:rounding}

실수 $x$를 $\mathbb{F}$의 원소로 바꾸는 사상을 $\fl:\mathbb{R}\to\mathbb{F}$라 쓰자. IEEE~754의 기본 규칙은 $x$에 가장 가까운 원소를 고르고, 두 원소가 똑같이 가까우면 가수가 짝수인 쪽을 고르는 \term{가장 가까운 짝수로의 반올림}{round to nearest, ties to even}이다.\footnote{동점에서 늘 위로 올리면 반올림 오차의 평균이 0이 되지 않아, 긴 계산에서 오차가 한쪽으로 쌓인다. 짝수 쪽을 고르면 동점의 절반은 올리고 절반은 내리게 되어 이 치우침이 사라진다.}

\begin{theorem}[반올림 오차의 상대 한계]
\label{thm:rounding-model}
$\beta=2$이고 $x$가 정규수의 범위 $2^{e_{\min}}\le|x|<2^{e_{\max}+1}$에 있으면
\[
\fl(x)=x(1+\delta),\qquad |\delta|\le u:=2^{-p}
\]
이다. $u$를 \term{반올림 단위}{unit roundoff}라 한다.
\end{theorem}

\begin{proof}
$x>0$이라 해도 일반성을 잃지 않는다. $2^{e}\le x<2^{e+1}$인 정수 $e$를 잡으면, $[2^{e},2^{e+1}]$ 안의 $\mathbb{F}$의 원소는 간격 $2^{e-p+1}$로 놓여 있다. 가장 가까운 원소까지의 거리는 간격의 절반을 넘지 않으므로
\[
|\fl(x)-x|\le\tfrac12\cdot2^{e-p+1}=2^{e-p}\le 2^{-p}x
\]
이다. 마지막 부등호는 $2^{e}\le x$에서 나온다. $\delta=(\fl(x)-x)/x$로 두면 된다.
\end{proof}

\begin{figure}[t]
\centering
\begin{tikzpicture}
\begin{axis}[sdplot, width=6.6cm, height=4.6cm,
  xmode=log, log basis x=2, xmin=0.5, xmax=8, ymin=0, ymax=0.14,
  xtick={0.5,1,2,4,8}, xticklabels={$\tfrac12$,$1$,$2$,$4$,$8$}, ytick={0,0.0625,0.125},
  yticklabels={$0$,$u/2$,$u$},
  xlabel={$x$}, ylabel={$|\fl(x)-x|/x$}]
\addplot+[no marks] table[x=x,y=relerr]{figures/data/fl-relerr-toy.dat};
\addplot+[no marks, domain=0.5:8] {0.125};
\end{axis}
\end{tikzpicture}
\caption{\cref{fig:float-line}의 형식($p=3$, $u=2^{-3}$)에서 반올림의 상대오차. 실선은 $|\fl(x)-x|/x$, 파선은 \cref{thm:rounding-model}의 한계 $u$이다. 오차는 원소 위에서 $0$이고 두 원소의 한가운데에서 가장 크며, 가로축을 로그 눈금으로 그리면 이진 구간마다 같은 무늬가 되풀이되고, 한계에 가장 가까이 가는 것은 각 구간의 왼쪽 끝 근처이다.}
\label{fig:relerr}
\end{figure}

\cref{thm:rounding-model}의 가치는(\cref{fig:relerr}) 오차를 \emph{상대적으로} 잰다는 데 있다. $x$가 $10^{-300}$이든 $10^{300}$이든, binary64는 그것을 유효숫자 약 16자리로 담는다. 한편 이 정리는 $x$를 담는 단계의 오차만 말해 준다는 점도 기억해 두자. 예컨대 $0.1$은 이진수로 유한하게 적히지 않으므로 컴퓨터 안의 \code{0.1}은 이미 참값과 약 $5.55\times10^{-18}$만큼 다르고, 그래서 \code{0.1 + 0.2}는 \code{0.30000000000000004}가 된다.

\begin{aside}[패트리어트 미사일과 0.1초]
1991년 2월 25일, 걸프 전쟁 중 사우디아라비아 다란에 떨어진 스커드 미사일을 미군의 패트리어트 방공 체계가 요격하지 못해 병사 28명이 목숨을 잃었다. 미국 회계감사원의 조사에 따르면 원인은 시간 계산이었다. 체계는 0.1초마다 늘어나는 정수 계수기에 $0.1$을 곱해 시각을 얻었는데, $0.1$을 24비트 고정소수점 레지스터에 담으며 소수 23자리에서 잘라 버렸다. 한 번의 곱마다 약 $9.54\times10^{-8}$초가 모자랐고, 체계를 100시간 동안 켜 둔 끝에 그 오차는 약 0.343초로 불어나 있었다. 초속 1.6\,km가 넘는 스커드에게 0.343초는 500\,m가 넘는 거리이다.
\end{aside}

\subsection{뺄셈이 정확할 때}

반올림은 대개 오차를 낳지만, 연산의 결과가 이미 $\mathbb{F}$에 있으면 오차가 없다. 다음 보조정리는 가까운 두 수의 뺄셈이 언제나 그렇다는 것을 말한다. 뒤에서 오차 없는 변환을 만들 때 핵심이 되는 사실이다.

\begin{lemma}[Sterbenz]
\label{lem:sterbenz}
$x,y\in\mathbb{F}$이고 $y/2\le x\le 2y$이면 $x-y\in\mathbb{F}$이다. 따라서 $\fl(x-y)=x-y$이다.
\end{lemma}

\begin{proof}
가정에서 $x,y\ge0$이다. $x\ge y$인 경우만 보이면 충분하다. 그렇지 않으면 두 수를 바꾸고 부호를 뒤집으면 된다. $x=m_x2^{e_x-p+1}$, $y=m_y2^{e_y-p+1}$로 쓰자. $y\le x$이므로 $e_y\le e_x$이고, 따라서
\[
x-y=\bigl(m_x2^{e_x-e_y}-m_y\bigr)\,2^{e_y-p+1}
\]
에서 괄호 안은 음이 아닌 정수이다. 한편 $x\le2y$에서 $x-y\le y$이므로 괄호 안의 정수는 $m_y<2^{p}$를 넘지 않는다. 가수가 $2^{p}$보다 작은 정수이고 지수가 $e_y$이므로 $x-y\in\mathbb{F}$이다.
\end{proof}

\section{합을 정확하게}
\label{sec:summation}

$n$개의 수를 왼쪽부터 차례로 더하면, 앞에서 생긴 반올림 오차가 뒤의 모든 덧셈을 거쳐 전해지므로 오차의 한계가 $n$에 비례해 자란다. 1965년 William Kahan은 각 덧셈에서 잃어버린 아랫자리를 따로 모아 두었다가 다음 덧셈에 되돌려 주는 간단한 방법으로 이 한계를 $n$과 거의 무관하게 만들었다(\cref{alg:kahan}).

\begin{algorithm}[t]
\caption{Kahan의 보정 합}
\label{alg:kahan}
\begin{algorithmic}[1]
\Require $x_1,\ldots,x_n\in\mathbb{F}$
\Ensure $s\approx\sum_{i=1}^{n}x_i$
\State $s\gets0$, $c\gets0$ \Comment{$c$: 지금까지 잃어버린 아랫자리}
\For{$i\gets1$ \textbf{to} $n$}
  \State $y\gets x_i-c$ \Comment{지난번에 잃은 것을 되돌린다}
  \State $t\gets s+y$ \Comment{$y$의 아랫자리가 여기서 잘린다}
  \State $c\gets(t-s)-y$ \Comment{잘린 양을 꺼낸다}
  \State $s\gets t$
\EndFor
\State \Return $s$
\end{algorithmic}
\end{algorithm}

\cref{alg:kahan}의 다섯째 줄 $c\gets(t-s)-y$는 수학적으로는 언제나 $0$이다. 이 식이 의미를 갖는 것은 오직 $t$가 반올림된 값이기 때문이며, 컴파일러가 대수 법칙을 써서 이 식을 $0$으로 ``최적화''해 버리면 알고리즘은 평범한 합으로 돌아간다. 이 미묘함은 \cref{lem:sterbenz}와 함께 다음 장에서 오차 없는 변환을 다룰 때 다시 만난다.

\begin{instardust}[\code{csrc/base/pairwise.hpp}]
stardust의 모든 부동소수점 합은 Kahan의 방법 대신 \emph{쌍마다 더하기}(pairwise summation)를 쓴다. 수열을 반으로 나누어 각각을 더한 뒤 둘을 더하는 일을 되풀이하면, 각 항이 거치는 덧셈의 수가 $\log_2 n$ 정도로 줄어 오차의 한계도 $\log n$에 비례하게 된다. 나누는 자리를 데이터의 길이만으로 정하므로, 스레드를 몇 개 쓰든 덧셈의 순서가 같고 결과는 비트까지 같다. 길이 512 이하의 조각은 독립된 누산기 16개로 더하는데, 이 수가 어디서 나왔는지는 \cref{ch:float}의 범위를 넘어 컴퓨터 구조를 다루는 장에서 설명한다.
\end{instardust}

\begin{remark}
쌍마다 더하기가 Kahan의 방법보다 오차 한계는 조금 크지만 더 많이 쓰이는 까닭은 속도이다. Kahan의 방법은 덧셈 하나에 연산 넷을 쓰고, 그 넷이 앞의 결과를 기다리는 사슬을 이룬다.
\end{remark}

